\section{Conclusion}\label{sec:conclusion}

We tested whether honesty steering works by suppressing vocabulary or by changing downstream computation. Our token-conditional projection zeros a steering vector's calibration-averaged first-order direct contribution to a chosen token set. A coherence-gated operating point and a random-subspace baseline make the resulting depth statistic interpretable.

CAA shifts honesty-coded vocabulary without a statistically significant positive gain at its selected coherent setting; benchmark-only selection instead chooses a collapsed setting. Mass-mean improves truthfulness within the coherence gate, and most of its effect survives aligned-token readout excision. At selected operating points across four models, CAA produces no statistically significant positive gain. Mass-mean succeeds on three, and most of each successful gain survives aligned-64 excision.

Useful steering is model-dependent, but the observed truthfulness gains are predominantly downstream under the tested readout. Researchers should treat a benchmark gain as a hypothesis, not as mechanistic evidence: verify coherent generation and held-out behavior before interpreting depth, then compare targeted excision with an equal-rank random control.
